Die Formelsammlung gibt die Dichte bei \TAgO, $\rho_0 = \rhGlO$ (FoTaBe S.~187), und $\gamma = \gaGlO$ (FoTaBe S.~191). Beim Erwärmen um $\Delta\vartheta = \dTgF = \dTg$ bleibt die Masse gleich, das Volumen wächst auf $V = V_0\left(1+\gamma\,\Delta\vartheta\right)$. Damit ist
\begin{align}
 \rho &= \frac{m}{V} = \frac{m}{V_0\left(1+\gamma\,\Delta\vartheta\right)} = \rhHF \\
   &= \frac{\rhGl}{1+\gaGl\cdot\dTg} \\
   &= \rhH \approx \result{\rhHP{0}{3}}
\end{align}
